
\section{Related Works}
\label{section2}


\subsection{Multi-Round Language Model Generation}
It has been widely noted that with generative language models, it is possible to get the desired answer through multiple rounds of interaction. Based on this core idea, several types of multi-round generation methods have been proposed to solve a wide range of problems. These methods include: 


\textbf{(1) Chain-of-Thought}. \citet{wei2022chain} proposes chain-of-thought (CoT), a prompt strategy that guides the LLM to generate regularized intermediate reasoning steps that lead from the initial question to the final answer. CoT has been proven to enable LLM{\color{black}s} to solve complex problems including dynamic programming, whereas non-CoT prompts fail to do so \citep{feng2023towards}. Empirically, variants of CoT have been proposed, including faithful CoT that tries to avoid LLMs lying in the reasoning stage \citep{lyu2023faithfulchainofthoughtreasoning}, multimodal CoT that enables CoT for vision context{\color{black}s} \citep{zhang2024multimodalchainofthoughtreasoninglanguage}, and tree-of-thought (ToT) that builds up a complex reasoning tree \citep{yao2023treethoughtsdeliberateproblem}.
%\citet{li2024planninglargelanguagemodels} empirically shows that conversational LLM agents augmented with CoT prompts better follow the task requirements. %\dcheng{rephrase: Yet, it is unclear how the CoT empowers/reduces the convergence round of conversation.}

\textbf{(2) Self-correction and Self-refinement.} LLMs' self-correction and refinement abilities have received significant attention recently \citep{madaan2024self, shinn2023reflexion, kim2023language}. Through explicit multi-round reflection, self-correction iteratively improves the accuracy, reliability, and robustness of LLM outputs, especially in complex reasoning tasks where the CoT reasoning process might be flawed \citep{madaan2024self}, by eliminating hallucinations~\citep{liu2024intrinsicselfcorrectioncapabilityllms}. %\dcheng{Rephrase: From the theoretical perspective,} self-correction abilities of LLMs are attributed to two properties of autoregressive generation. \citet{liu2024large} proves that self-correction is an implicit self-verification structure \citep{gero2023selfverification,weng2023large} that reduces the inherent hallucination during autoregressive generation. On the other hand, \citet{liu2024intrinsicselfcorrectioncapabilityllms, liu2024intrinsicselfcorrectionenhancedmorality} argue through empirical observations that model uncertainty of LLMs decreases and converges to some value after rounds of self-correction. %\dcheng{Rephrase: The convergence of explicit constraint from human prompts empirically leads to convergence of LLM answers, but a theoretical analysis on whether the LLMs themselves could reach \textit{any} answer is still missing.}

\textbf{(3) Multi-Agent Debating}.
Apart from improving a single-agent LLM system, another line focuses on the collaboration of multiple LLM agents \citep{li2023camel,wei2023multipartychatconversationalagents,hao2023chatllmnetworkbrainsintelligence}. Findings suggest that multi-agent collaboration through both debating or even majority voting can often outperform single-agent LLM with explicit constraints \citep{huang2024large,wu2023autogenenablingnextgenllm}. \citet{pmlr-v235-khan24a} further shows that debating with more persuasive LLMs results in better answers. % While multiple works have showcased the successful utilization of multi-agent LLM systems, the theoretical analyses behind such success are relatively inadequate. \dcheng{check if this is relevant--} Before LLMs, in the field of game theory, reinforcement learning of multi-agent systems has been shown to reach Nash equilibrium on both two-player imperfect information games \citep{perolat2020poincarerecurrenceconvergenceimperfect} and two-player zero-sum perfect information games \citep{silver2017masteringchessshogiselfplay}. An early work from \citet{irving2018aisafetydebate} proves that the complexity class analog of the debate can answer any question in PSPACE using only polynomial time judges. %\dcheng{Rephrase: Yet, under the context of large language models and limited rounds, how debating between LLM agents could foster implicit constraints on the convergence of the discussed topic remains unknown.}

Overall, the essence of these methods is to impose internal or external interventions on multi-round sequence generation, which is an optimization constraint in the generalization process.

%these methods provide internal or external interventions during the multi-round generation process to place constraints on the error propagation of the sequence generation process.

\subsection{Theoretical Analysis on Language Models}

\textbf{Approximation Ability of Transformers.} The approximation capabilities of Transformer architectures have been extensively studied in recent years. \citet{yuntransformers} establishes that Transformers are universal approximations of continuous permutation equivariant sequence-to-sequence functions. \cite{wei2022statistically} introduces a statistically meaningful approximation, demonstrating that overparameterized neural networks can approximate Boolean circuits and Turing machines with generalization guarantees. \citet{schuurmans2023memory} proves that Transformers can simulate Turing machines in the presence of conditional memory. In terms of limitation{\color{black}s}, \citet{dong2021attention} demonstrates that pure attention mechanisms could suffer from rank collapse, leading to a loss of expressive power with increased network depth. \citet{hahn2020theoretical} analyzes the theoretical limitations of self-attention in neural sequence models, highlighting challenges in modeling hierarchical structures. %An interesting conclusion is that, 
\cite{caivocabulary} proves that the composition of words in a finite vocabulary can approximate any continuous function in a compact domain. Unlike their setting, we show that a finite context Transformer can approximate a Turing Machine by up to infinite rounds of generation.

{\color{black}
\textbf{Simulating Turing Machine with Neural Networks.} Numerous scholars have proposed using neural networks to simulate Turing machines to verify the computational power of neural networks~\citep{siegelmann1992computational,perez2021attention,perez2019turing,wei2022statistically,graves2014neural}. \cite{siegelmann1992computational} point out that recurrent neural networks (RNNs) of infinite precision can make simulations of Turing machines, implying the Turing-completeness of RNNs. \cite{chung2021turing} further implements a finite precision RNN simulation of Turing Machine using an external memory module. However this Turing-complete property cannot be directly inherited to a finite-window Transformer, because the recurrent learning model in RNNs essentially aggregates information from time 0. While the Transformer model via self-attention does not pass on the information, but rather, it is in the form of KV cache~\citep{gemodel}, which results in information beyond the context window being encoded only on newly generated tokens, a completely different computational paradigm from RNNs. For the Turing completeness of Transformer, \cite{bhattamishra2020computational} prove that finite-precision Transformers are \textbf{not} Turing-complete. \cite{perez2021attention} prove that only hard-attention can be Turing complete. But Transformers still can approximate the simulation of Turing Machine within a certain margin of error~\citep{wei2022statistically}.
Yet, these studies have not examined the relationship between approximation precision and quantization precision. We reveal in this work the approximation of a Turing machine that can be reached through multiple rounds of Transformer's generation, as well as a high-level relationship between the precision of Transformer's numerical quantization and the approximation error tolerance.
}

\textbf{Learnability and Sample Complexity.} Probably Approximately Correct learning is a framework in computational learning theory that provides a formal way to understand how well an algorithm can learn a target function from a set of examples. The concept was introduced by \citet{valiant1984theory} and helped analyze how quickly and accurately a learning algorithm can generalize from data. For sequence modeling, \citet{schuurmans1995sequential} proposes a kind of sequential PAC learning where online stopping rules are used to minimize the number of samples required for PAC learning. \citet{gavalda2006pac} studies the PAC learnability of hidden Markov models. \citet{pmlr-v235-malach24a} extends this framework to the prediction of the next token, but there are assumptions of infinite context-awareness, as well as ignoring error propagation. So far, to the best of our knowledge, there are no previous works that study the PAC learnability and sample complexity of long sequence generation with a context window-limited auto-regressive generative model. Making up for the shortcomings of the above works, in this paper, we study the PAC learnability of sequence generation and apply it to a multi-round generation task to theoretically qualitatively study the effect of multi-round generation on the sample complexity of the learning of {\color{black}the} long sequence generation task.


